% généré par scripts/06_textes.py
\dosmeta{Année de classement : 367 (directe : nom consulaire lu). Références : ICUR I 940 = ILCV 2878. Rattachement selon l’EDB : \emph{Incertae originis} (le volume I des ICVR réunit les inscriptions dont la provenance est perdue). Support : Plaque de marbre. Inscription gravée. Reproductions signalées par l’EDB : dessin, ICVR I (de Rossi) (droits réservés ; non reproduites). Fiche EDB : Antonio Enrico Felle, première version 2014. Notice du catalogue p.~\pageref{not:110}.}
\txtbloc{Texte transmis (EDCS, état 2022)}{\lnum{1}ἐν εἰρήνῃ ἀνεπάη Βεναινᾶτος ἀπ(ὸ) κο(μήτων)\par
\lnum{2}ζήσας ἔτη οεʹ ὃς κῖτε ἐν τῷδε τόπῳ ἐξελτὼν\par
\lnum{3}ἐκ τοῦ βίου τῇ πρ(ὸ) δʹ ἰδ(ῶ)ν μαρ(τίων) ὑπ(ατείᾳ) Λουππικί(νου)\par
\lnum{4}καὶ Ἰοβίνου Θεοφίλα ἐποίησα συμβίῳ\par
\lnum{5}Benenato Theofila fecit uxor eius}
\txtbloc{Texte de l’EDB (EDB19681)}{\lnum{1}ἐν εἰρήνῃ ἀνεπάη Βεναινᾶτος ἀπ(ὸ) κο(μήτων)\par
\lnum{2}ζήσας ἔτη οε´ ὃς κῖτε ἐν τῷδε τόπῳ ἐξελτὼν\par
\lnum{3}ἐκ τοῦ βίου τῇ π̂ρ(ὸ) δ´ ἰδ(ῶ)ν μαρ(τίων) ὑπ(ατείᾳ) Λουππικί(νου)\par
\lnum{4}καὶ ᾿Ιοβίνου Θεοφίλα ἐποίησα συμβίῳ\par
\lnum{5}Benenato Theofila fecit uxor eius}
